% Borrador focal de adición, NO demostración de la igualdad de dos vías.
% Puede incorporarse su proposición positiva sin alterar fórmulas actuales.
% La nota PROCEDENCIA_Y_ALCANCE.md delimita la compatibilidad no reconstruida.
\paragraph{Exact dependence on the truncation boundary.}
\label{ampl-alfa-frontera-exacta}
The regional coordinates and the dodecaphasic register come from the
common APP--TRIT--TPK generation. Once produced, their positional
composition has an explicit dependence on readout
depth. Let $B=1000$,
\[
 Z_j=P_j+E_j-\Phi_j-K_j^{\mathrm{per}},\qquad
 y_N=\sum_{j=1}^{N}Z_jB^{-j},\qquad
 R_{N+1}(Z)=\sum_{r\geq0}Z_{N+1+r}B^{-r-1}.
\]
The complete series in \eqref{eq:alpha-serie} satisfies exactly
\[
 \alpha_A=y_N+B^{-N}R_{N+1}(Z),\qquad
 -2<R_{N+1}(Z)<2.
\]
This identity combines the contribution of the window and that of its
extension, both expressed in the same coordinates.

\begin{proposition}[Boundary transport and canonical cylinder]
\label{ampl-alfa-prop-frontera}
Let $A_j^{(N,t)}$ be the output of finite positional division with
$c_{N+1}=t\in\mathcal C=\{-2,-1,0,1,2\}$, and let
$a_N^{(t)}=\sum_{j=1}^{N}A_j^{(N,t)}B^{-j}$. For the registers in the
domain considered, we have
\begin{equation}
 \alpha_A-a_N^{(t)}
 =B^{-N}\bigl(R_{N+1}(Z)-t\bigr),
 \qquad
 |\alpha_A-a_N^{(t)}|<4B^{-N}.
 \label{ampl-alfa-eq-frontera}
\end{equation}
The quotient of the complete extension is
$c_{N+1}^{\infty}=\lfloor R_{N+1}(Z)\rfloor$. With this boundary,
\[
 0\leq\alpha_A-a_N^{(c_{N+1}^{\infty})}<B^{-N},
\]
so the resulting window is the canonical prefix of order $N$.
\end{proposition}
\begin{proof}
Telescoping \eqref{eq:alpha-identidad-local} gives
$a_N^{(t)}=y_N+tB^{-N}-c_1$. Since $Z_1=7$, every incoming state
belongs to $\mathcal C$, and the first dividend lies between $5$ and $9$,
we have $c_1=0$. Substituting this equality into the series
decomposition proves \eqref{ampl-alfa-eq-frontera}. The bound follows from
$R_{N+1}(Z)\in(-2,2)$ and $t\in\mathcal C$.

For complete canonical normalization,
$c_{N+1}^{\infty}=R_{N+1}(Z)-R_{N+1}(A)$ and
$R_{N+1}(A)\in[0,1)$. Integrality of the quotient implies
$c_{N+1}^{\infty}=\lfloor R_{N+1}(Z)\rfloor$. The error formula
then becomes
$\alpha_A-a_N^{(c_{N+1}^{\infty})}=B^{-N}R_{N+1}(A)$, proving
membership in the canonical cylinder.
\end{proof}

The proposition identifies the meaning of the right boundary: it is the
integer part of the signed tail expressed at the scale of the next
position. Calculation over the five boundaries preserves this information
as a finite collection of possibilities until the blocks to the
left stabilize. The quantitative contraction comes from depth and
tail transport; its object is the normalization of
coordinates already generated.

\paragraph{Graded domain of the analytic resolvent.}
\label{ampl-alfa-dominio-resolvente}
The analytic composition displayed in
\S\ref{subsec:alpha-resolvente} acts on
$\operatorname{span}\{e_7,e_8,e_9\}$ and satisfies $Ne_9=0$. Its
evaluation produces exactly $T_{7:9}$; degree $9$ is the last
component of that domain. The identity $N^3=0$ explains why the
resolvent is expressed through three summands. Describing an
extension to higher degrees additionally requires specification of the
transport between successive graded domains and its action on the
initial component of each domain. Local nilpotence and global
compatibility thus serve distinct mathematical functions.
